\chapter{Conclusões}

Este trabalho apresentou uma abordagem integrada de AM e Pesquisa Operacional PO para apoiar o processo de montagem de elencos no futebol profissional. 
Inicialmente, foi desenvolvido um modelo preditivo capaz de estimar o desempenho esperado de equipes a partir das características individuais dos jogadores. 
Os resultados obtidos indicam que o modelo proposto foi capaz de explicar aproximadamente 55\% da variabilidade da pontuação final dos clubes, evidenciando 
sua capacidade de capturar padrões relevantes nos dados e produzir estimativas consistentes de desempenho coletivo.

A partir dessas previsões, o modelo foi utilizado como função objetivo em um problema de otimização voltado à seleção de jogadores sob restrições orçamentárias. 
Para explorar o espaço de soluções, foram implementadas duas heurísticas: um algoritmo guloso baseado em substituições iterativas e um método de busca local.

Do ponto de vista prático, os resultados obtidos indicam que a abordagem proposta pode contribuir para tornar o processo de tomada de decisão na montagem de 
elencos mais estruturado e orientado por dados. No contexto do futebol brasileiro, onde decisões de contratação frequentemente são influenciadas por fatores 
subjetivos ou avaliações intuitivas, ferramentas baseadas em análise quantitativa podem oferecer suporte adicional à identificação de jogadores que se encaixem 
de forma mais eficiente em determinadas estruturas de equipe, respeitando restrições financeiras.

Assim, este trabalho demonstra que a integração entre técnicas de AM e métodos de PO representa um caminho promissor para apoiar processos de recrutamento e 
planejamento esportivo. Ao transformar dados de desempenho em estimativas de impacto coletivo e utilizá-las em modelos de decisão, torna-se possível reduzir 
a dependência exclusiva da intuição na escolha de atletas e avançar em direção a abordagens mais analíticas na gestão do futebol. Nesse sentido, 
os resultados apresentados reforçam o potencial da integração entre Aprendizado de Máquina e Pesquisa Operacional como uma abordagem promissora para 
transformar dados de desempenho esportivo em ferramentas quantitativas de apoio à tomada de decisão, contribuindo para o avanço de métodos analíticos 
aplicados à gestão e ao planejamento no futebol profissional.

\section{Trabalhos Futuros}

Como continuidade desta pesquisa, algumas direções podem ser exploradas em trabalhos futuros. Uma possibilidade consiste no aprimoramento dos modelos de AM, 
buscando capturar de forma ainda mais precisa a relação entre o desempenho individual dos jogadores e o desempenho coletivo das equipes. 
A incorporação de novas variáveis explicativas, como características táticas das equipes, métricas de interação entre jogadores ou informações contextuais das partidas, 
pode contribuir para aumentar a capacidade preditiva dos modelos.

Outra linha de investigação envolve o desenvolvimento de métodos de otimização mais avançados para o problema de montagem de elencos, incluindo 
metaheurísticas ou abordagens híbridas capazes de explorar o espaço de soluções de maneira mais eficiente. Além disso, futuras pesquisas podem considerar 
restrições adicionais relevantes ao contexto real do futebol profissional, como limitações contratuais, regras de competições e requisitos específicos de 
formação de elenco.

Por fim, a ampliação da base de dados utilizada, incorporando informações de outras ligas e temporadas, pode contribuir para tornar os modelos mais 
robustos e generalizáveis, ampliando o potencial de aplicação prática da abordagem proposta.




